\chapter{Generator liczb losowych w technice cyfrowej}
\label{ch:generator-cyfrowy}

Celem niniejszego rozdziału jest przedstawienie podstaw teoretycznych oraz stanu
wiedzy niezbędnych do zrozumienia budowy i zasady działania generatora liczb
losowych zaimplementowanego w układzie FPGA. Struktura rozdziału odpowiada
logicznemu następstwu zagadnień: terminologia i klasyfikacja generatorów, model
funkcjonalny generatora fizycznego wraz z wymaganiami, jakie musi on spełniać,
a następnie koncepcja fizycznej realizacji losowości w strukturach cyfrowych.
Szczególną uwagę poświęcono roli oscylatora pierścieniowego jako elementu
wprowadzającego niedeterministyczność, mechanizmowi powstawania jitteru,
architekturom ekstrakcji entropii oraz zjawiskom, które tę entropię ograniczają.

Aparat pomiarowy służący do ilościowej oceny tak zdefiniowanego generatora,
to jest miary entropii oraz metody analizy statystycznej strumienia bitów,
przedstawiono w rozdziale kolejnym. Podział ten jest celowy: niniejszy rozdział
odpowiada na pytanie, \emph{skąd} bierze się losowość i jakie zjawiska mogą ją
zaburzyć, natomiast rozdział następny na pytanie, \emph{jak} zaburzenia te
wykryć i zmierzyć.

\section{Generatory liczb losowych - klasyfikacja, model funkcjonalny i wymagania}
\label{sec:klasyfikacja-wymagania}

\subsection{Rola liczb losowych}
\label{sec:rola-liczb-losowych}

Liczby losowe są punktem wejściowym dla znacznej części mechanizmów
bezpieczeństwa wielu systemów informatycznych. Wykorzystywane są do generowania
kluczy kryptograficznych, wektorów inicjalizujących czy wartości jednorazowych
(nonce). Poza zagadnieniami z dziedziny kryptografii losowość jest podstawą
metod Monte Carlo, algorytmów probabilistycznych, generowania danych testowych
oraz gier losowych.

Kluczowym zagadnieniem w tych zastosowaniach jest błąd w generatorze. Nie jest
to awaria funkcjonalna systemu, lecz cicha utrata bezpieczeństwa. System,
którego generator ma zbyt małą entropię, działa poprawnie i przechodzi testy
funkcjonalne, a jednocześnie produkuje klucze możliwe do odtworzenia przez
atakującego. Praktycznym przykładem tego zjawiska jest sytuacja udokumentowana
w badaniu Heningera i in.~\cite{heninger2012mining}, w którym wśród kilku
milionów kluczy publicznych zebranych z urządzeń sieciowych wykryto liczne
klucze powtarzające się lub współdzielące czynniki, co przypisywano
niedostatecznej entropii w chwili pierwszego uruchomienia urządzenia.

Wniosek jest następujący: jakość generatora nie może być wnioskowana wyłącznie
z poprawnego działania systemu, lecz powinna być przedmiotem osobnego,
ilościowego badania statystycznego. To założenie stanowi punkt wyjściowy
niniejszej pracy.

\subsection{Klasyfikacja generatorów}
\label{sec:klasyfikacja}

W literaturze oraz w dokumentach normalizacyjnych~\cite{nist_sp80090c}
przyjmuje się następujący podział.

Generatory deterministyczne DRNG (ang. \emph{Deterministic Random Number
Generator}), potocznie zwane również PRNG, wytwarzają ciąg bitów jako
deterministyczną funkcję stanu wewnętrznego zainicjowanego ziarnem (seed).
Ciąg jest w pełni odtwarzalny, a entropia wyjścia, niezależnie od jego
długości, nie przekracza entropii ziarna. Zaletami takiego rozwiązania są
wysoka przepustowość, powtarzalność (istotna np. przy symulacjach) oraz
przewidywalny koszt sprzętowy. Wadą jest zależność bezpieczeństwa od ziarna,
które musi pochodzić z fizycznego źródła.

Generatory fizyczne, prawdziwie losowe (TRNG, ang. \emph{True Random Number
Generator}), pozyskują nieprzewidywalność z niepowtarzalnych procesów
fizycznych: szumu termicznego, zjawisk kwantowych, metastabilności elementów
bistabilnych. Ciąg wyjściowy nie jest odtwarzalny, a entropia rośnie wraz
z długością ciągu. Wadami takiego rozwiązania są ograniczona i zależna od
warunków pracy przepustowość, wrażliwość na temperaturę, napięcie zasilania
i zakłócenia oraz konieczność stałego monitorowania poprawności działania.

Generatory hybrydowe łączą źródło fizyczne z przetwarzaniem algorytmicznym:
albo źródło fizyczne zasila generator deterministyczny (rozwiązanie dominujące
w praktyce przemysłowej~\cite{killmann2011ais31}), albo surowe bity są
poddawane kondycjonowaniu kryptograficznemu. Uzyskuje się w ten sposób wysoką
przepustowość przy zachowaniu nieprzewidywalności.

Generator badany w niniejszej pracy należy do drugiej kategorii. Analizie
podlega strumień bitów pobierany bezpośrednio ze źródła fizycznego, przed
kondycjonowaniem kryptograficznym, ponieważ dopiero taki strumień pozwala
ocenić jakość samego źródła entropii. Zastosowanie silnej funkcji
kondycjonującej maskuje bowiem defekty źródła, nie usuwając ich.

\subsection{Model funkcjonalny generatora fizycznego}
\label{sec:model-funkcjonalny}

Zarówno AIS-31~\cite{killmann2011ais31}, jak i SP 800-90B~\cite{turan2018sp80090b}
opisują generator fizyczny jako łańcuch przetwarzania, w którym każdy blok
podlega odrębnej ocenie:

\begin{enumerate}
  \item \textbf{Źródło entropii} - element analogowy lub analogowo-cyfrowy,
        w którym powstaje nieprzewidywalność (w układzie FPGA: szum fazowy
        oscylatora, metastabilność).
  \item \textbf{Digitalizacja (próbkowanie)} - przetworzenie zjawiska
        ciągłego na ciąg bitów. Wynik nazywany jest \emph{surowym ciągiem
        losowym} (ang. \emph{raw random numbers}).
  \item \textbf{Kondycjonowanie / korekcja algorytmiczna} - poprawa
        właściwości statystycznych i kompresja entropii (blok opcjonalny).
  \item \textbf{Dodatkowe testy wbudowane} - test całkowitej awarii
        (\emph{total failure test}), testy startowe oraz testy bieżące
        (\emph{health tests}, \emph{online tests}), których zadaniem jest
        wykrycie degradacji źródła w czasie pracy.
  \item \textbf{Interfejs wyjściowy} - buforowanie i udostępnianie
        wewnętrznych liczb losowych aplikacji.
\end{enumerate}

Podział ten ma bezpośrednie znaczenie dla części doświadczalnej pracy.
Właściwości statystyczne mierzone na wyjściu bloku drugiego opisują jakość
samego źródła, natomiast pomiar wykonany po bloku trzecim opisuje już
właściwości łańcucha przetwarzania jako całości. Rozdzielenie tych dwóch
punktów pomiarowych pozwala stwierdzić, czy poprawa parametrów statystycznych
wynika ze wzrostu entropii źródła, czy wyłącznie z działania korekcji
algorytmicznej.

\subsection{Wymagania stawiane generatorowi w układzie FPGA}
\label{sec:wymagania}

Praktyczne wymagania stawiane generatorowi w układzie FPGA można ująć
w następujących kategoriach:

\begin{itemize}
  \item \textbf{entropia na bit} - ilość entropii wewnętrznych liczb losowych
        przypadająca na bit wyjściowy. Dokumenty normalizacyjne posługują się
        przy tym różnymi miarami: AIS-31~\cite{killmann2011ais31} przyjmuje
        entropię Shannona, natomiast SP 800-90B~\cite{turan2018sp80090b}
        min-entropię, czyli miarę najgorszego przypadku. Ponieważ min-entropia
        jest miarą konserwatywną i bezpośrednio odpowiada odporności na
        odgadnięcie wartości, w niniejszej pracy traktuje się ją jako
        właściwe kryterium oceny; szczegółowe porównanie obu miar wraz
        z uzasadnieniem tego wyboru przedstawiono w rozdziale następnym.
  \item \textbf{przepustowość} - od kilkudziesięciu kbit/s (generowanie
        kluczy) do setek Mbit/s (maskowanie, protokoły wysokoprzepustowe).
  \item \textbf{zasoby} - liczba komórek LUT, bloków pamięci
        i zasobów specjalnych (PLL/MMCM), a także pobór mocy.
  \item \textbf{odporność środowiskowa} - zachowanie parametrów w pełnym
        zakresie temperatury, napięcia i procesu (PVT).
  \item \textbf{odporność na ataki} - pasywne (analiza emisji
        elektromagnetycznej, analiza poboru mocy) i aktywne (wstrzykiwanie
        częstotliwości, manipulacja zasilaniem).
  \item \textbf{przenośność} - możliwość implementacji w różnych rodzinach
        układów bez utraty jakości.
  \item \textbf{testowalność} - możliwość obserwacji surowego strumienia
        bitów przed kondycjonowaniem, będąca warunkiem koniecznym zarówno
        certyfikacji, jak i badania przeprowadzonego w niniejszej pracy.
\end{itemize}

Wymagania te pozostają ze sobą w konflikcie. Zwiększanie przepustowości przez
skrócenie okresu próbkowania obniża entropię na bit, ponieważ zmniejsza czas
akumulacji jitteru. Zwielokrotnienie liczby źródeł podnosi entropię, lecz
zwiększa zużycie zasobów i pobór mocy, a przy zbyt ciasnym rozmieszczeniu
prowadzi do zjawisk opisanych w podrozdziale~\ref{sec:zjawiska-ograniczajace}.
Projektowanie generatora jest zatem poszukiwaniem kompromisu, którego
poprawność można ocenić jedynie pomiarowo.

\section{Oscylator pierścieniowy jako źródło entropii}
\label{sec:oscylator-pierscieniowy}

\subsection{Zasada działania}
\label{sec:zasada-dzialania}

Oscylator pierścieniowy (ang. \emph{ring oscillator}) jest jednym z
najprostszych układów generujących sygnał okresowy w technologiach
cyfrowych. Składa się z nieparzystej liczby inwerterów połączonych
szeregowo w zamkniętą pętlę, gdzie wyjście ostatniego elementu jest
podłączone do wejścia pierwszego.

Zasada działania jest intuicyjna: sygnał logiczny krążący w pętli jest
kolejno odwracany przez każdy inwerter. Ponieważ liczba inwerterów jest
nieparzysta, układ nie może osiągnąć stabilnego stanu logicznego.
Zamiast tego powstaje ciągła oscylacja pomiędzy stanem wysokim i niskim.
Częstotliwość tych oscylacji zależy głównie od sumy opóźnień propagacji
poszczególnych bramek. Dla pierścienia złożonego z $N$ inwerterów o średnim
opóźnieniu propagacji $t_p$ okres oscylacji wynosi w przybliżeniu

\begin{equation}
  T_{osc} \approx 2 N t_p ,
  \label{eq:okres-oscylatora}
\end{equation}

co oznacza, że częstotliwością można sterować liczbą elementów w pętli.
Zależność ta jest jednak wyłącznie przybliżeniem wartości średniej; istotne
z punktu widzenia generowania losowości jest to, że rzeczywisty okres
poszczególnych cykli podlega fluktuacjom.

\subsection{Jitter jako zjawisko fizyczne i jego model stochastyczny}
\label{sec:jitter}

Istotną cechą oscylatora pierścieniowego jest jego podatność na zjawiska
fizyczne zachodzące w strukturze krzemowej. Szumy termiczne, fluktuacje
napięcia zasilania, zmiany temperatury oraz lokalne różnice technologiczne
powodują losowe zmiany opóźnień propagacji. Skutkiem tego jest
nieregularność okresu sygnału, określana jako jitter, stanowiąca praktyczne
źródło entropii w wielu konstrukcjach
TRNG~\cite{sunar2006provably,saarinen2021entropy}.

Dla oceny generatora kluczowe jest rozróżnienie dwóch składowych tego
zjawiska, ponieważ tylko jedna z nich niesie entropię:

\begin{itemize}
  \item \textbf{jitter losowy} (ang. \emph{random jitter}) - wynika ze
        zjawisk o charakterze szumowym, przede wszystkim szumu termicznego
        w kanałach tranzystorów. Fluktuacje kolejnych okresów można
        traktować jako niezależne, o rozkładzie w przybliżeniu normalnym.
        Jest to jedyna składowa, którą wolno uznać za źródło entropii.
  \item \textbf{jitter deterministyczny} (ang. \emph{deterministic jitter}) -
        wynika ze sprzężeń z sygnałami obecnymi w układzie: zegarem
        systemowym, przebiegami na liniach zasilania, aktywnością sąsiednich
        bloków logicznych. Ma charakter okresowy lub skorelowany
        z aktywnością układu, jest zatem z zasady przewidywalny i nie
        stanowi entropii, mimo że powiększa obserwowaną wariancję okresu.
\end{itemize}

Rozróżnienie to jest podstawą modelu stochastycznego wymaganego przez
AIS-31~\cite{killmann2011ais31}. Przyjmując, że fluktuacje kolejnych okresów
są niezależne i mają odchylenie standardowe $\sigma_1$, wariancja
akumulowanej fazy po czasie obserwacji $t_{acc}$ narasta liniowo z liczbą
przebytych okresów:

\begin{equation}
  \sigma^2(t_{acc}) = \frac{t_{acc}}{T_{osc}}\,\sigma_1^{2},
  \label{eq:akumulacja-jitteru}
\end{equation}

skąd odchylenie standardowe fazy rośnie proporcjonalnie do
$\sqrt{t_{acc}}$. Konsekwencja praktyczna tej zależności jest zasadnicza:
entropia pojedynczego bitu zależy od stosunku zakumulowanej niepewności
fazy do okresu sygnału próbkowanego. Jeżeli okres próbkowania jest zbyt
krótki, niepewność nie osiąga skali umożliwiającej nieprzewidywalny wynik
kwantyzacji, a kolejne bity pozostają wzajemnie zależne. Zależność między
parametrami próbkowania a entropią na bit stanowi przedmiot modeli
przedstawionych w literaturze~\cite{ma2014entropy,saarinen2021entropy},
przy czym wspólnym wnioskiem jest istnienie dolnej granicy okresu
próbkowania, poniżej której generator traci właściwości statystyczne
niezależnie od pozostałych elementów konstrukcji.

Model ten wyjaśnia zarazem, dlaczego badanie generatora nie może ograniczać
się do rozkładu zer i jedynek. Niedostateczna akumulacja jitteru objawia się
przede wszystkim jako zależność między kolejnymi bitami, natomiast jitter
deterministyczny jako składowa okresowa w strumieniu. Obie te postacie
defektu są niewidoczne dla pomiaru samej częstości bitów, a wykrywalne
metodami przedstawionymi w rozdziale następnym.

\subsection{Architektury ekstrakcji entropii}
\label{sec:architektury-ekstrakcji}

Sam oscylator pierścieniowy nie jest jeszcze generatorem. Entropia zawarta
w jitterze musi zostać przeniesiona na ciąg bitów, a sposób realizacji tego
przeniesienia decyduje o entropii na bit oraz o przepustowości. W literaturze
wyodrębnia się kilka rozwiązań podstawowych.

\textbf{Próbkowanie elementarne (ERO, ang. \emph{elementary ring oscillator}).}
Sygnał oscylatora o wysokiej częstotliwości jest próbkowany przerzutnikiem
taktowanym sygnałem o częstotliwości znacznie niższej, pochodzącym
z niezależnego źródła. Rozwiązanie wymaga minimalnych zasobów, lecz osiągnięcie
zadowalającej entropii na bit wymusza bardzo niską częstotliwość próbkowania,
ponieważ konieczna jest akumulacja jitteru zgodnie
z zależnością~\eqref{eq:akumulacja-jitteru}. Przepustowość jest zatem
niewielka.

\textbf{Wiele oscylatorów sumowanych funkcją XOR (MURO, ang. \emph{multiple
ring oscillators}).} Rozwiązanie polega na zwielokrotnieniu liczby
niezależnych oscylatorów i zsumowaniu ich spróbkowanych wyjść funkcją XOR.
Podejście to, przedstawione w pracy Sunara i in.~\cite{sunar2006provably},
pozwala zwiększyć przepustowość przy zachowaniu entropii, a funkcja XOR
zapewnia dodatkowo, że pojedynczy oscylator o obniżonej jakości nie
zdominuje wyniku. Warunkiem poprawności tej konstrukcji jest jednak
\emph{wzajemna niezależność} oscylatorów, założenie w praktyce trudne do
spełnienia, o czym mowa w podrozdziale~\ref{sec:zjawiska-ograniczajace}.
Właśnie ta architektura, jako wykorzystująca zwielokrotnione źródła,
stanowi punkt odniesienia dla generatora badanego w niniejszej pracy.

\textbf{Próbkowanie koherentne (COSO, ang. \emph{coherent sampling}).}
Wykorzystuje się dwa oscylatory o bardzo zbliżonych, lecz nieidentycznych
częstotliwościach i mierzy liczbę okresów jednego z nich mieszczącą się
w okresie zdudnienia obu sygnałów. Wynik pomiaru jest tym bardziej czuły na
jitter, im mniejsza jest różnica częstotliwości, co pozwala uzyskać wysoką
entropię przy umiarkowanym zużyciu zasobów. Rozwiązania tego typu są
podstawą konstrukcji przenośnych, zachowujących właściwości niezależnie od
rodziny układu~\cite{peetermans2019portable}.

\textbf{Konstrukcje oparte na pętlach PLL oraz na sterowanym starcie
i zatrzymaniu oscylatora.} Zamiast swobodnie oscylującego pierścienia
stosuje się generowanie sygnałów o precyzyjnie ustalonym stosunku
częstotliwości albo cykliczne uruchamianie i zatrzymywanie oscylatora, tak
aby faza początkowa każdego cyklu pomiarowego była niezależna od
poprzedniego. Rozwiązania tej klasy pozwalają zwiększyć przepustowość przy
zachowaniu deterministycznej latencji i przenośności między
platformami~\cite{matuszewski2024startstop}.

Wybór architektury jest zatem wyborem punktu na krzywej kompromisu między
entropią na bit, przepustowością i zużyciem zasobów. Żadna z nich nie
gwarantuje jednak samoczynnie jakości strumienia, ponieważ wszystkie opierają
się na założeniach dotyczących niezależności źródeł i skali jitteru, których
spełnienie w konkretnej implementacji podlega weryfikacji pomiarowej.

\subsection{Zjawiska ograniczające entropię}
\label{sec:zjawiska-ograniczajace}

Realna entropia generatora opartego na oscylatorach jest z reguły niższa od
wartości wynikającej z modelu idealnego. Odpowiadają za to następujące
zjawiska.

\textbf{Wzajemna synchronizacja oscylatorów.} Oscylatory pierścieniowe
umieszczone blisko siebie w strukturze układu oddziałują poprzez wspólną sieć
zasilania oraz sprzężenia pojemnościowe między ścieżkami. Skutkiem jest
tendencja do wzajemnego przyciągania faz i częstotliwości, prowadząca do
częściowej lub pełnej synchronizacji. Zjawisko to jest szczególnie groźne
w architekturach zwielokrotnionych: jeżeli spośród $n$ zaimplementowanych
oscylatorów efektywnie niezależnych pozostaje jedynie $k < n$, to entropia
strumienia odpowiada $k$, a nie $n$ źródłom, przy niezmienionym zużyciu
zasobów. Co istotne, synchronizacja nie musi objawiać się zmianą rozkładu
zer i jedynek na wyjściu; ujawnia się natomiast jako korelacja między
strumieniami pochodzącymi z poszczególnych źródeł oraz jako struktura
okresowa w strumieniu wypadkowym.

\textbf{Dominacja jitteru deterministycznego.} Jeżeli udział składowej
deterministycznej w obserwowanej wariancji okresu jest znaczny, oszacowanie
entropii oparte na całkowitej wariancji jest zawyżone. Praktyczną
konsekwencją jest zawyżona ocena jakości generatora przy jednoczesnej
obecności składowej periodycznej w strumieniu.

\textbf{Wrażliwość na warunki pracy.} Zarówno częstotliwość oscylacji, jak
i amplituda jitteru zależą od temperatury i napięcia zasilania. Generator
poprawny w warunkach nominalnych może utracić właściwości statystyczne
w skrajnych punktach zakresu pracy, co uzasadnia wymóg testów bieżących
w modelu funkcjonalnym z podrozdziału~\ref{sec:model-funkcjonalny}.

\textbf{Podatność na ataki aktywne.} Wstrzykiwanie sygnału o częstotliwości
zbliżonej do częstotliwości oscylatora albo celowa modulacja napięcia
zasilania mogą wymusić synchronizację źródeł, a w skrajnym przypadku
sprowadzić strumień wyjściowy do postaci przewidywalnej. Odporność na tego
rodzaju oddziaływania jest jednym z motywów stosowania architektur
zwielokrotnionych~\cite{sunar2006provably}.

Wspólną cechą powyższych zjawisk jest to, że żadne z nich nie objawia się
w sposób oczywisty w działaniu układu. Ich wykrycie wymaga celowanego
pomiaru: analizy zależności czasowych w strumieniu oraz porównania strumieni
pochodzących z poszczególnych źródeł. Wyznacza to zakres badania
przeprowadzonego w dalszej części pracy.

\section{Implementacja generatora w strukturach FPGA}
\label{sec:implementacja-fpga}

Układy FPGA (ang. \emph{Field Programmable Gate Array}) stanowią
elastyczną platformę do implementacji cyfrowych generatorów liczb
losowych, w tym konstrukcji bazujących na jitterze oscylatorów
pierścieniowych. Architektura oparta na programowalnych blokach logicznych
i rekonfigurowalnych połączeniach umożliwia bezpośrednią realizację
struktur sprzętowych i ich równoległe
zwielokrotnienie~\cite{peetermans2019portable}.

W przeciwieństwie do rozwiązań programowych, implementacja generatora
w strukturze FPGA pozwala na kontrolę nad topologią, sposobem próbkowania
oraz logiką post-processingu. Jest to istotne, ponieważ zarówno jakość
źródła entropii, jak i mechanizm jej ekstrakcji wpływają na końcową
losowość strumienia bitów~\cite{ma2014entropy}.

Platforma ta wprowadza jednak ograniczenia nieobecne w realizacjach
dedykowanych. Rozmieszczenie elementów w strukturze układu jest wynikiem
działania narzędzi implementacyjnych, przez co odległości między
oscylatorami, a zatem i skala ich wzajemnego oddziaływania, nie są w pełni
kontrolowane przez projektanta. Narzędzia syntezy dążą ponadto do
eliminowania struktur kombinacyjnych zawierających pętle sprzężenia
zwrotnego, traktując je jako błąd projektowy. Zachowanie oscylatorów
w postaci zamierzonej wymaga zatem jawnego blokowania optymalizacji, co
sprawia, że jakość uzyskanego źródła entropii zależy nie tylko od koncepcji
układowej, lecz również od sposobu przeprowadzenia implementacji.

W praktyce platformy FPGA są często wybierane w zastosowaniach
wymagających wysokiej przepustowości i deterministycznej latencji,
a także tam, gdzie istotna jest przenośność rozwiązania pomiędzy
różnymi rodzinami układów. Przykładem są współczesne konstrukcje TRNG
projektowane tak, aby zachowywać właściwości statystyczne niezależnie
od producenta i szczegółów technologii~\cite{matuszewski2024startstop}.

\section{Podsumowanie}
\label{sec:podsumowanie-generator}

Z przedstawionego materiału wynikają trzy wnioski wyznaczające zakres części
doświadczalnej pracy.

Po pierwsze, entropia generatora fizycznego jest właściwością źródła, nie
strumienia bitów, i nie może być wnioskowana z poprawnego działania układu.
Jej ocena wymaga osobnego badania ilościowego, prowadzonego na surowym
strumieniu, przed kondycjonowaniem algorytmicznym.

Po drugie, entropia oscylatora pierścieniowego pochodzi wyłącznie ze
składowej losowej jitteru i zależy od czasu jego akumulacji względem okresu
próbkowania. Zarówno zbyt krótki okres próbkowania, jak i dominacja jitteru
deterministycznego prowadzą do obniżenia entropii, przy czym oba defekty
objawiają się jako zależności czasowe w strumieniu, nie jako zaburzenie
rozkładu zer i jedynek.

Po trzecie, architektury oparte na zwielokrotnionych źródłach opierają się
na założeniu ich wzajemnej niezależności, którego spełnienie w konkretnej
implementacji nie jest zagwarantowane ze względu na zjawisko wzajemnej
synchronizacji oscylatorów.

Wnioski te wskazują, jakie własności strumienia należy poddać pomiarowi:
rozkład brzegowy bitów, zależności czasowe w obrębie jednego strumienia
oraz zależności między strumieniami pochodzącymi z różnych źródeł. Miary
entropii pozwalające wyrazić te własności liczbowo oraz metody statystyczne
służące ich wykrywaniu przedstawiono w rozdziale następnym.